\section{GPT-1 与 GPT-2：从 Fine-Tuning 到 Zero-Shot Interface}

GPT 系列的早期演进不仅是参数增长，也在改变 task interface。GPT-1 采用通用预训练后 supervised fine-tuning；GPT-2 更强调任务可以被写进自然语言上下文，由同一个 language model zero-shot 完成。

\subsection{GPT-1：Pretrain then fine-tune}

\term{fine-tuning}（微调）是在预训练参数上继续用目标任务数据更新。GPT-1 用 decoder-style Transformer 进行 generative pretraining，再为分类、entailment、question answering 等任务构造输入格式并微调。

\lecturefigure{slide-14-gpt1-architecture.jpg}{GPT-1 将通用语言预训练与任务微调连接起来。}{官方视频 14:11；Radford et al. 2018。}

\begin{knowledgebox}{读图：架构复用不等于完全零改动}
左侧是统一 Transformer representation，右侧用 task-specific input transformation 和 output head 适配不同任务。Pretraining 提供初始化，fine-tuning 仍需要标签、超参数与每任务 checkpoint。它比从零训练高效，却没有形成统一 prompt-only interface。
\end{knowledgebox}

\subsection{GPT-2：Zero-shot reading comprehension}

本节先用 reading comprehension 说明统一文本接口。\term{zero-shot} 指目标任务上不提供训练 example 或参数更新，只通过自然语言上下文要求模型执行。GPT-2 将 question、context 和 answer continuation 拼成序列，观察 language modeling 是否隐式完成 QA。

\lecturefigure{slide-15-gpt2-reading-comprehension.jpg}{GPT-2 通过文本格式尝试 zero-shot reading comprehension。}{官方视频 15:21；Radford et al. 2019。}

\begin{knowledgebox}{读图：任务被转写为 continuation}
模型没有专用 QA head，而是预测 prompt 后的文本。评价必须定义 answer extraction、exact match 或 semantic equivalence；若 prompt format 变化大，结果可能显著波动。Zero-shot 证明 interface 的可塑性，不代表模型真正遵循任意 instruction。
\end{knowledgebox}

\subsection{Zero-shot summarization}

Summarization 同样被写成文档后接提示词与摘要。Language model 会利用训练数据中类似格式，产生较短 continuation。核心困难是 coverage、faithfulness 和 compression，不是语法流畅。

\lecturefigure{slide-16-gpt2-summarization.jpg}{GPT-2 将 summarization 重新表述为条件文本续写。}{官方视频 16:44。}

\begin{knowledgebox}{读图：摘要质量至少有三条轴}
Coverage 检查关键事实是否保留；faithfulness 检查是否引入原文没有的 claim；compression 检查长度与冗余。ROUGE 等 overlap metric 只覆盖部分维度。生成式模型可能写出更自然的摘要，同时更容易产生不可见事实错误。
\end{knowledgebox}

\subsection{Zero-shot translation}

接下来把同一 continuation 接口用于翻译。模型可通过“源句 + language marker”被诱导；GPT-2 的结果展示多任务 signal 可能从 web text 中被吸收，但与专门训练的 translation system 仍有明显差距。

\lecturefigure{slide-17-gpt2-translation.jpg}{GPT-2 zero-shot translation 展示 web pretraining 中的潜在跨语言 signal。}{官方视频 17:30。}

\begin{knowledgebox}{读图：比较 baseline 与任务定义}
图中的趋势应与 supervised translation baseline 对照。Language model 可能见过平行片段、模板或相似网站，因此 zero-shot 不等于“从未见过翻译关系”。应检查 contamination、language coverage 与低资源语言，而不是只看高资源平均分。
\end{knowledgebox}

\begin{warningbox}{Zero-shot 是 evaluation protocol，不是魔法属性}
是否 zero-shot 取决于目标任务是否参与训练、prompt 是否含示例、benchmark 是否泄漏、以及参数是否更新。Web-scale pretraining 很难证明“任务从未出现”，因此更准确的表述是目标 benchmark 上无显式 supervised fine-tuning。
\end{warningbox}

\subsection{本章小结}

GPT-1 通过 pretrain/fine-tune 复用表示，GPT-2 把任务转写为 continuation，向统一自然语言接口迈进。Zero-shot behavior 有价值，却高度依赖 prompt、data exposure 和 metric。

